\documentclass[10pt,a4paper]{amsart}
\usepackage[margin=19mm]{geometry}
\usepackage{amsmath,amssymb,mathtools}
\usepackage[scr=rsfs,cal=euler]{mathalfa}
\usepackage{xfrac}
\usepackage[unicode,hidelinks,bookmarks=false]{hyperref}
\newcommand{\ie}{i.e.\ }
\newcommand{\cf}{cf.\ }
\newcommand{\resp}{respectively\ }
\newcommand{\Subsection}{\subsection}
\providecommand{\nobreakdash}{\nobreak\hbox{-}}
\setlength{\emergencystretch}{2em}
\multlinegap0pt
\usepackage{mathtools}
\usepackage{thm-restate}
\newtheorem{theorem}{Theorem}[section]
\newtheorem{proposition}[theorem]{Proposition}
\newtheorem{corollary}[theorem]{Corollary}
\newtheorem{lemma}[theorem]{Lemma}
\newtheorem{remark}[theorem]{Remark}
\newtheorem{definition}[theorem]{Definition}
\newtheorem{example}[theorem]{Example}
\newcommand{\C}{\mathbb C}
\newcommand{\E}{\mathbb E}
\newcommand{\DN}{\mathcal D_N}
\newcommand{\St}{\mathrm{St}}
\newcommand{\op}{\mathrm{op}}
\newcommand{\norm}[1]{\left\lVert #1 \right\rVert}
\newcommand{\ip}[2]{\left\langle #1,#2\right\rangle}
\newcommand{\Arow}{\left\|\sum_{i=1}^n A_i^2\right\|_{\op}^{1/2}}
\DeclareMathOperator{\Tr}{Tr}
\DeclareMathOperator{\Rea}{Re}
\DeclareMathOperator{\Ima}{Im}
\DeclareMathOperator{\vecop}{vec}
\begin{document}
\title[A Sudakov--Fernique proof of Lehner-type edge bounds for mat]{A Sudakov--Fernique proof of Lehner-type edge bounds for matrix-valued GUE sums}
\author{Benoit Collins; Yuta Yamagishi}
\date{}
\maketitle
\noindent\emph{Source and licence.} A Sudakov--Fernique proof of Lehner-type edge bounds for matrix-valued GUE sums. Version arXiv:2606.21137v2 (CC BY 4.0). This material is distributed under the Creative Commons Attribution 4.0 International licence, http://creativecommons.org/licenses/by/4.0/, as recorded in the arXiv OAI licence metadata for this exact version and shown on the abstract page https://arxiv.org/abs/2606.21137v2. Full rights notice: licenses/arxiv-2606.21137-CC-BY-4.0.txt.\par\smallskip
\noindent\textit{Editorial scope.} Counted unit: the original \texttt{theorem} environment \texttt{thm:lehner} at source lines 157--176 together with its complete proof at lines 178--194; the delivered document keeps source lines 107--195 so that every referenced label and every citation resolves inside it. Native editable source is the paper's own concatenated TeX; the AMS wrapper is editorial formatting only. No publisher class, upstream script or unrelated artwork is redistributed.
\par\smallskip\noindent\textit{Reference note.} This article ships no compiled bibliography (biblatex); the entries delivered here are composed from the author's own BibTeX (.bib) fields, with no field invented and no entry dropped.
\section{Introduction}

Free probability was introduced by Voiculescu in the 1980s as a noncommutative probability theory in which the notion of independence is replaced by a noncommutative counterpart called freeness.  Although this theory was initially developed in the context of operator algebras, it has found many applications.
One of the first striking connections was with random matrix theory, when
Voiculescu showed that independent GUE matrices become asymptotically free in the limit as their size tends to infinity~\cite{Voiculescu1991}.
Strong asymptotic freeness strengthens this by asserting that the norms of polynomials in independent random matrices converge to the norms of the corresponding free operators in the reduced free product $C^*$-algebra.  This was first proved for GUE matrices by Haagerup and Thorbj\o rnsen~\cite{HaagerupThorbjornsen}, and quantitative forms were later developed by Collins, Guionnet and Parraud~\cite{CollinsParraudGuionnet}.

One common route from polynomial norm convergence to spectral estimates is the
linearization trick, which replaces a polynomial by a larger self-adjoint linear
pencil with matrix coefficients.  However, the strong-convergence proofs
themselves use several different mechanisms.  Haagerup and Thorbj\o rnsen's
original proof uses a master equation for the Stieltjes transform of the
spectrum; Collins, Guionnet and Parraud's approach relies on interpolating
between the GUE and a free product; other techniques use resolvent and
linearization methods~\cite{Anderson2013}, matrix
concentration~\cite{BBvH2023}, and Markov inequalities and interpolation
techniques~\cite{CGVTvHClassicalI,CGVvHClassical,vanHandelSurvey2025}.

The goal of the present paper is different.  We study directly a linear
matrix-coefficient model and show how far one can obtain non-asymptotic
spectral-edge bounds using only Gaussian concentration and comparison
inequalities.  This model is natural in its own right: it is a block Gaussian
matrix with a macroscopic covariance pattern.  The resulting finite-dimensional
error is not optimized, but in the regime of uniformly bounded positive
coefficients and $N=o(M)$ it has the correct free leading constant.  We do not
prove the corresponding asymptotic lower bound.  A subsequent preprint of
Stojnic~\cite{Stojnic2026} claims the analogous asymptotic equality for a real
symmetric Gaussian Kronecker model with positive semidefinite coefficients; the
main setting of that work takes the coefficient dimensions to be fixed.  The
same leading-constant estimate applies to
certain signed families whose covariance map has a positive semidefinite Kraus
decomposition.


Specifically, we focus on the linear matrix-coefficient model
\[
	H_M=A_0\otimes I_M+\frac1{\sqrt M}\sum_{i=1}^n A_i\otimes G_i,
\]
where $A_0,A_1,\ldots,A_n\in M_N(\C)$ are Hermitian.  In the
free limit, $M^{-1/2}G_i$ is replaced by a free semicircular variable
$s_i=l_i+l_i^*$, where $l_i$ are the canonical creation operators on full Fock
space.  The following edge formulas are extracted from Lehner's semicircular
operator-norm formula.

The coefficient family determines the completely positive covariance map
\begin{equation}
	\label{eq:covariance-map}
	\eta(X)=\sum_{i=1}^n A_iXA_i,\qquad X\in M_N(\C).
\end{equation}


\begin{theorem}[Edge formulas from Lehner's semicircular formula]
	\label{thm:lehner}
	Let $l_1,\ldots,l_n$ be the canonical creation operators on full Fock space and
	let $a_0,a_1,\ldots,a_n\in M_m(\C)$ with $a_0=a_0^*$ and $a_i=a_i^*$.  Put
	\[
		x=a_0\otimes I+\sum_{i=1}^n a_i\otimes(l_i+l_i^*).
	\]
	Then
	\[
		\lambda_{\max}(x)
		=
		\inf_{z\succ0}\lambda_{\max}\left(z+a_0+\sum_{i=1}^n a_iz^{-1}a_i\right)
	\]
	and
	\[
		\lambda_{\min}(x)
		=
		\sup_{z\prec0}\lambda_{\min}\left(z+a_0+\sum_{i=1}^n a_iz^{-1}a_i\right).
	\]
\end{theorem}

\begin{proof}
	Here \(\lambda_{\max}(x)=\sup\operatorname{sp}(x)\) and
	\(\lambda_{\min}(x)=\inf\operatorname{sp}(x)\).
	Lehner states an operator-norm formula under the additional assumption
	\(a_0\succeq0\) and \(n\ge2\) in \cite[Corollary~1.5]{Lehner}.  The case
	\(n=1\) follows by adjoining a zero coefficient and an additional free
	semicircular variable.  For arbitrary Hermitian \(a_0\), add a sufficiently
	large scalar multiple \(cI\) so that both \(x+cI\) and the matrices in
	Lehner's variational formula are positive.  The operator norm is then the
	right edge; subtracting \(c\) yields the first formula.  See also
	\cite[Lemma~2.4]{BBvH2023}, whose \(\varepsilon=+1\) term gives the
	corresponding right-edge variational expression after the change of variables
	\(Z\leftrightarrow Z^{-1}\), and Parmaksiz--van
	Handel~\cite{ParmaksizVanHandel2025} for the corresponding edge and
	singular-value formulation.  The left-edge formula follows by applying the
	right-edge formula to \(-x\).
\end{proof}


\begin{thebibliography}{99}
\bibitem[]{Anderson2013} Anderson, Greg W.. Convergence of the largest singular value of a polynomial in independent {W}igner matrices. The Annals of Probability 41 (2013) 2103--2181
\bibitem[]{BBvH2023} Bandeira, Afonso S.; Boedihardjo, March T.; van Handel, Ramon. Matrix concentration inequalities and free probability. Inventiones Mathematicae 234 (2023) 419--487
\bibitem[]{CGVTvHClassicalI} Chen, Chi-Fang; Garza-Vargas, Jorge; Tropp, Joel A.; van Handel, Ramon. A new approach to strong convergence. Annals of Mathematics 203 (2026) 555--602
\bibitem[]{CGVvHClassical} Chen, Chi-Fang; Garza-Vargas, Jorge; van Handel, Ramon. A new approach to strong convergence {II}. {The} classical ensembles. Geometric and Functional Analysis (2026)
\bibitem[]{CollinsParraudGuionnet} Collins, Beno{\^i}t; Guionnet, Alice; Parraud, Felix. On the operator norm of non-commutative polynomials in deterministic matrices and iid {GUE} matrices. Cambridge Journal of Mathematics 10 (2022) 195--260
\bibitem[]{HaagerupThorbjornsen} Haagerup, Uffe; Thorbj{\o}rnsen, Steen. A new application of random matrices: {$\mathrm{Ext}(C^*_{\mathrm{red}}(F_2))$} is not a group. Annals of Mathematics 162 (2005) 711--775
\bibitem[]{Lehner} Lehner, Franz. Computing norms of free operators with matrix coefficients. American Journal of Mathematics 121 (1999) 453--486
\bibitem[]{ParmaksizVanHandel2025} Parmaksiz, Emre; van Handel, Ramon. Computing extreme singular values of free operators. arXiv preprint arXiv:2510.23987 (2025)
\bibitem[]{Stojnic2026} Stojnic, Mihailo. An {RDT} based confirmation of {Lehner}'s formula for {Kronecker}--{Gaussian} matrices. arXiv preprint arXiv:2607.26551 (2026)
\bibitem[]{Voiculescu1991} Voiculescu, Dan. Limit laws for random matrices and free products. Inventiones Mathematicae 104 (1991) 201--220
\bibitem[]{vanHandelSurvey2025} van Handel, Ramon. The strong convergence phenomenon. Current Developments in Mathematics 2025 (2026) 177--261
\end{thebibliography}
\end{document}
